\section{Alignment geometry and computational parameters}\label{sec:geometry}
This section fixes the input conventions and the interval states used
throughout the algorithm. It also records the geometric inequalities that
allow child intervals to be combined even when their target endpoints do
not agree.

\subsection{Paths, states and connection costs}
Index string boundaries starting at zero. An alignment path has steps
$(1,0),(0,1),(1,1)$, of cost one, one, and the indicator that the two
characters on a diagonal step differ. Such a path gives an edit sequence.
Conversely, the original characters that survive an edit sequence occur in
the final string in their original order. Deleting other original
characters, inserting the remaining output characters and changing each
surviving character whose label differs gives an alignment of no greater
cost. Thus minimum alignment cost equals edit distance.

For an interval $I=(l,u]$ of the first string, an interval state is
$(I,q)$, where $q=(q^-,q^+)$ is an integer pair with
$0\le q^-\le q^+\le |y|$. When $I$ is fixed, we also call $q$ a state. Set
\[
 c(I,q)=\ED\bigl(x_{(l,u]},y_{(q^-,q^+]}\bigr),\qquad
 z(I,q)=(q^+-q^-)-(u-l).
\]
An endpoint's \emph{shift} is its target index minus its first-string index.
In particular, $z$ is the change in shift between the two endpoints.
For pairs use $|q-r|_1=|q^--r^-|+|q^+-r^+|$.

\begin{lemma}\label{lem:geometry}
For every interval state, $c(I,q)\ge |z(I,q)|$. Both $c(I,\cdot)$ and
$|z(I,\cdot)|$ are $1$-Lipschitz in the endpoint metric.
If $I$ is partitioned, in order, into $I_1,\ldots,I_m$, and
$\alpha=(s_1,\ldots,s_m)$ is any list of valid child states, define
\[
 \conn_r(\alpha)=|s_1^--r^-|
   +\sum_{i=1}^{m-1}|s_{i+1}^--s_i^+|+|s_m^+-r^+|.
\]
Then
\begin{align}
 c(I,r)&\le\sum_i c(I_i,s_i)+\conn_r(\alpha),\label{eq:connection}\\
 |z(I,r)|&\le\sum_i|z(I_i,s_i)|+\conn_r(\alpha).
 \label{eq:shift-connection}
\end{align}
The connection inequalities also hold for finite target intervals in any
extension of $y$ beyond its stored endpoints.
\end{lemma}
\begin{proof}
Only insertions and deletions change length, proving $c\ge|z|$.
Changing a target substring's endpoints costs at most $|q-r|_1$:
remove or insert endpoint portions if the substrings overlap, and delete
and insert the whole substrings if they are disjoint. The triangle
inequality gives the claim for $c$; the claim for $|z|$ follows from its
formula.

Edit each child to its specified target substring and concatenate the
results. Consider the walk on the target's boundary line that starts at
$r^-$, jumps to $s_1^-$, traverses each child target forward, makes the
successive connecting jumps, and ends at $r^+$. Expand each jump into unit
steps. From this walk choose, in order, one forward crossing of every unit
interval between $r^-$ and $r^+$. Successive first arrivals at the next
boundary give such a choice. The letters on these selected crossings,
in their walk order, form exactly the parent's target substring. Delete
the letters on unselected child crossings, and insert the letters on
selected connecting crossings. By net displacement, the number of
unselected forward crossings equals the number of backward crossings.
All backward crossings belong to connecting jumps, since child targets
are traversed forward. Thus deletions cost at most the total backward
jump length, while insertions cost at most the total forward jump length.
Their sum is $\conn_r(\alpha)$, proving~\eqref{eq:connection}.
This argument uses only the letters at finitely many target positions,
so it also proves the assertion for extensions of $y$.
Finally, sum the signed shift changes of child intervals and connecting
jumps and take absolute values to obtain~\eqref{eq:shift-connection}.
\end{proof}

Figure~\ref{fig:connection-cost} illustrates why the connection term permits
both gaps and overlaps between consecutive child targets.
\begin{figure}[htbp]
\centering
\begin{tikzpicture}[
  x=0.72cm,y=0.84cm,>=Latex,
  interval/.style={->,very thick},
  connection/.style={->,densely dashed,thick},
  endpoint/.style={circle,fill,inner sep=1.6pt},
  every node/.style={font=\small}]
\draw[->] (-0.35,1.1) -- (10.6,1.1);
\node at (5,1.95) {target boundary};
\foreach \t in {0,1,3,4,6,7,9,10} {
  \draw (\t,1.04) -- (\t,1.16);
  \node[above] at (\t,1.16) {$\t$};
}
\node[left] at (-0.4,0) {$I_1$};
\node[left] at (-0.4,-1.4) {$I_2$};
\node[left] at (-0.4,-2.8) {$I_3$};
\draw[interval] (1,0) -- (4,0)
  node[midway,above] {$s_1=(1,4)$};
\draw[interval] (3,-1.4) -- (6,-1.4)
  node[midway,below] {$s_2=(3,6)$};
\draw[interval] (7,-2.8) -- (9,-2.8)
  node[midway,below] {$s_3=(7,9)$};
\foreach \x/\y in {1/0,4/0,3/-1.4,6/-1.4,7/-2.8,9/-2.8}
  \node[endpoint] at (\x,\y) {};
\node[endpoint,label=above left:{$r^-=0$}] (start) at (0,0.45) {};
\node[endpoint,label=below right:{$r^+=10$}] (finish) at (10,-3.25) {};
\draw[connection] (start) -- (1,0) node[midway,below left] {$1$};
\draw[connection] (4,0) -- (3,-1.4) node[midway,left] {$1$};
\draw[connection] (6,-1.4) -- (7,-2.8) node[midway,right] {$1$};
\draw[connection] (9,-2.8) -- (finish) node[midway,above right] {$1$};
\node[align=center] at (5,-4.05)
  {$\conn_r(\alpha)=|1-0|+|3-4|+|7-6|+|10-9|=4$};
\end{tikzpicture}
\caption{Connection charges for three consecutive source children.
Solid arrows show their prescribed target intervals; dashed arrows join
these intervals in child order, with the outside endpoints $r=(0,10)$.
The jump from $4$ to $3$ points backward, so this walk is not itself an
alignment path. Lemma~\ref{lem:geometry} gives
$c(I,r)\le\sum_{i=1}^3c(I_i,s_i)+4$.
Only the horizontal target-coordinate displacement of a dashed arrow is
charged; the vertical separation of rows is for display.}
\label{fig:connection-cost}
\end{figure}


\subsection{Uniform access to the accuracy input}
The rational accuracy is supplied in a fixed canonical binary encoding of
its positive numerator and denominator in lowest terms, on a separate
read-only input.
For example, encode a positive integer by its binary length in unary,
a delimiter, and its binary digits, and concatenate the two encodings.
The input is promised to specify $0<\varepsilon<1$. Prefix inspection is
ordinary input access; no numerical operation on the rational is free.

The algorithm first handles equality, empty strings and $|x|\le1$
exactly. These cases require $O(N+1)$ string work. On other inputs define
$H$ as the least power of two at least $\lceil\log_2(N+2)\rceil$, and
$\ell$ as the least power of two at least $\max(1,\log_2 H)$.
Inspect at most $H^2+O(1)$ bits of the accuracy input. If its complete
encoding has not been obtained, use exact alignment dynamic programming.
The unread suffix is neither copied nor addressed. All addresses actually
used for this prefix have $O(\log(N+2))$ bits. If the encoding is obtained,
its integers have $O(H^2)$ bits, and exact arithmetic decides the test
$\ell\varepsilon\ge10$ with polynomial work in $H$.

Use exact dynamic programming also if $\ell<2^{1000}$ or
$\ell\varepsilon<10$. It uses $O((N+1)^2)$ word operations and can store
only two rows. Together with the bounded parameter inspection, this
charges every operation in these branches. For each fixed rational
encoding, the inspection bound eventually exceeds its length, and both
numeric tests eventually pass as $N$ grows. Thus the exact branches affect
only a bounded range of $N$ for each fixed $\varepsilon$. They preserve
the uniform algorithm and its approximation guarantee on that range.
No assertion about efficient joint dependence on encoding length or
$1/\varepsilon$ is needed.

\subsection{Parameters and the fixed interval tree}
All remaining arguments concern the large-input branch. Logarithms with
subscript two are binary, and $\ln$ denotes the natural logarithm.
The coarse estimator and interval tree use the branching factors $B$ and
$M$ below. The factor $A_{\rm init}$ is the initial seed accuracy, $F$
controls its reduction across passes, and $P$ controls endpoint-grid
resolution:
\begin{align*}
 p_*&=\left\lfloor(\log_2\ell)^{1/4}\right\rfloor,
 &B&=H^{p_*},&A_{\rm init}&=H^{p_*+4},\\
 M&=2^{\lfloor(\log_2 H)/20\rfloor},
 &P&=H^{p_*+20},&F&=\ell.
\end{align*}
The counts $T$ and $S$ govern preliminary propagation and refinement
iterations, respectively. The slopes $L_0$ and $L$ penalize changes in
target endpoints:
\[
 T=2^{\lceil(3/4)\log_2H\rceil},\qquad S=H/\ell^2,
 \qquad L_0=32F,\qquad L=128F.
\]
Finally, fix the error tolerances and the sampling and optimization
parameters. Their roles are specified at the corresponding constructions:
\begin{align*}
 \delta&=\ell^{-10},&\kappa&=\tau=\ell^{-4},
 &\theta&=\ell^{-40},\\
 \gamma&=\ell^{-20},&\eta&=\ell^5,&Q_1&=\ell^{80}.
\end{align*}
The essential separation is between $H=\Theta(\log(N+2))$ and
$\ell=\Theta(\log H)$: the branching factors and grid resolutions are
powers of $H$, while the tolerances are powers of $\ell$. The inequalities
below quantify the slack available between these scales.

Every integer parameter here is computable using bit lengths and integer
tests. For example, $p_*$ is the largest integer whose fourth power is
at most the integer $\log_2\ell$. The displayed fractional powers and
logarithms require no real-arithmetic oracle. All parameters used as dyadic
grid spacings are powers of two, and $S$ is a positive integer in this
branch.

We will repeatedly use the generous bounds
\begin{equation}\label{eq:parameter-slack}
 H\ge2^{\ell/2},\qquad \ell^{200}\le H^{1/100},\qquad
 p_*+21\le\sqrt\ell.
\end{equation}
The first follows from the definition of $\ell$; the other two follow at
$\ell\ge2^{1000}$ and become stronger as $\ell$ increases. More explicitly,
$200\log_2\ell\le\ell/200$ and
$(\log_2\ell)^{1/4}+21\le\sqrt\ell$ hold at this cutoff, and their
right sides thereafter grow faster. Factors polynomial in $\ell$ and
$p_*$ will consequently be much smaller than any fixed positive power
of $H$ needed below.

Partition $x$ into a complete $M$-ary tree of height
\[
 J=\left\lceil\log_M\max(1,|x|)\right\rceil.
\]
At a node $(l,u]$ use cuts
$l_i=l+\lfloor i(u-l)/M\rfloor$, $0\le i\le M$.
Empty intervals are allowed, and node labels distinguish repeated
intervals. Nodes need only be generated when requested. At depth $j$ the
width is at most $\lceil|x|/M^j\rceil$, and leaves have width at most one.
The parameters give
\begin{equation}\label{eq:tree-bounds}
 J\le100H/\ell,\qquad M^J\le M\max(1,|x|).
\end{equation}
Indeed $\log_2M\ge (\log_2H)/20-1$ and
$\log_2H\ge\ell/2$, whereas $\log_2|x|\le H$.
The second bound follows directly from the ceiling defining $J$.

At a leaf of width zero the exact cost is the target length. At a leaf
containing one character, the cost is one for an empty target, and otherwise
is the target length minus the indicator that this character occurs there.
Sort the pairs (symbol, position) of $y$ lexicographically by comparison
sorting. A binary search then tests occurrence in any target interval in
$O(\log(N+2))$ time, with $O(N\log(N+2))$ preprocessing work and linear
storage. The word-label assumption makes each comparison a bounded number
of word operations. Fresh padding symbols used later can be represented by
a tag together with a label; no search for an unused alphabet value is
needed.

All subsequent tables are indexed by this fixed tree and valid endpoint
pairs. The total domain is polynomial in $N+2$: there are at most
$2M^J$ nodes and at most $(|y|+1)^2$ pairs per node. A simultaneous
probability guarantee over this domain is meaningful even though the
algorithm will evaluate only a small part of it.
